\documentclass[12pt]{article}
\title{Generalized Linear Model}
\date{\today}
\usepackage{amsmath}

\setlength{\parindent}{0pt} % don't indent the line
\begin{document}
\maketitle

Family: binomial or Poisson. This specifies the variance function. This allows the generalized linear model to relax the assumption of constant variance of residuals.\\

Link functions: transform response variable (e.g. count, probability) by some functions (e.g. natural log, logit). The transformed response variable is called linear predictor $\eta$ (predictor here means predicted value): $\eta = \beta_0 + \beta_1x_1 + \cdots$ \\

$\eta$ has a linear relationship with $\beta_0 + \beta_1x_1 + \cdots$:
\[ \eta = \beta_0 + \beta_1x_1 + \cdots \]

Apply inverse link $f^{-1}(x)$ to $\eta$ to transform it back to y: 
\[ f^{-1}(\eta) = f^{-1}(\beta_0 + \beta_1x_1 + \cdots) \]

\[ y = f^{-1}(\beta_0 + \beta_1x_1 + \cdots) \] \\

The model will predict $\hat{y}$ by applying inverse link to $\beta_0 + \beta_1x_1 + \cdots$. \\

\section{Link function: natural log (when y is count)}
Link: $\eta = ln(y)$ \\

Inverse link: $ y = e^\eta$ \\

Example: \\

There is a generalized linear model where $\beta_0 = 1, \beta_1 = 1, x_1 = -5$:
\[ ln(y) = \eta = \beta_0 + \beta_1x_1\] \\

The predicted value of $ln(y)$ is $1 + 1 \times (-5) = -4$ (link scale). The predicted value of $y$ is $e^{-4} \approx 0.018$ (data scale).


\section{Link function: logit (when y is probability)}
odds = $\frac{p}{1-p}$ \\
log-odds (logit) = $ln(\frac{p}{1-p})$ \\

Link: $\eta = ln(\frac{y}{1-y})$ \\
Inverse link: $y = \frac{1}{1+e^{-\eta}}$ \\

Example: \\

Calculate the probability $p$ when log-odds is 0. \\

\begin{equation*}
\begin{split}
ln(\frac{p}{1-p}) &= 0 \\
e^{log(\frac{p}{1-p})} &= e^0 \\
\frac{p}{1-p} &= 1 \\
p &= 1-p \\
p &= 0.5
\end{split}
\end{equation*}

Calculate the probability $p$ when log-odds is 1. \\

\begin{equation*}
\begin{split}
ln(\frac{p}{1-p}) &= 1 \\
e^{ln(\frac{p}{1-p})} &= e^1 \\
\frac{p}{1-p} &= e \\
p &= e-ep \\
(1+e)p &= e \\
p &= \frac{e}{1+e} \approx 0.73
\end{split}
\end{equation*}




\end{document}